\documentclass[11pt]{article}

\usepackage[margin=1in]{geometry}
\usepackage[T1]{fontenc}
\usepackage[utf8]{inputenc}
\usepackage{amsmath}
\usepackage{amssymb}
\usepackage{booktabs}
\usepackage{enumitem}
\usepackage{microtype}
\usepackage[numbers,sort&compress]{natbib}
\usepackage{xcolor}
\usepackage{hyperref}

\hypersetup{
  colorlinks=true,
  linkcolor=blue!50!black,
  citecolor=blue!50!black,
  urlcolor=blue!50!black
}

\newcommand{\code}[1]{\texttt{#1}}
\newcommand{\artifact}[1]{\path{#1}}
\newcommand{\modelname}{CodeLeWM}

\title{\modelname{}: A Reproducible Code World-Model Study with
  Negative Action-Use Results and a Narrow Downstream Reranking Slice}
\author{
  Abdelhamid Bakhta\\
  StarkWare\\
  \url{https://github.com/AbdelStark/CodeLeWM}
}
\date{Draft of June 8, 2026}

\begin{document}
\maketitle

\begin{abstract}
Joint-embedding predictive architectures learn by predicting future latent
states rather than reconstructing observations. \modelname{} ports this recipe
to Python code transitions and asks whether a compact latent world model can
add value when ranking candidate code changes. The answer is mixed. The
original commit-edit substrate and the v0.2 action-swap intervention are
negative: text-action retrieval reaches Recall@1 0.263 and MRR 0.370, while
the no-action baseline reaches Recall@1 0.441 and MRR 0.533. Execution-trace
substrates repair collapse and produce healthy diagnostics: v0.9 retrieval
beats all retrieval controls, mutation surprise AUC is 1.0 on both seeds, and
the downstream HumanEval WS-D replay reaches CodeLeWM pass@1 0.9787 on both
seeds, beating no-action by +10.64 and +8.51 points with confidence intervals
excluding zero. The aggregate public claim remains closed. MBPP-Plus is
saturated with CodeLeWM, no-action, and lexical all at pass@1 1.0000; broad
semantic surprise has no scorable semantic pairs; representation probes do not
beat controls; and the final score rows do not expose a standalone \code{p\_pass}
score key. We therefore present \modelname{} as a reproducible negative and
diagnostic study with a narrow HumanEval WS-D positive slice, not as evidence
that the model generally improves coding.
\end{abstract}

\section{Introduction}

World models are attractive because they promise useful predictions without
directly imitating every observation. In joint-embedding predictive
architectures (JEPAs), the model predicts the latent of a future state rather
than reconstructing pixels or tokens
\citep{assran2023selfsupervisedlearningimagesjointembedding,bardes2024revisitingfeaturepredictionlearning}.
LeWorldModel applies the idea to action-conditioned visual dynamics from
pixels \citep{maes2026leworldmodelstableendtoendjointembedding}. \modelname{}
asks whether the same latent-transition framing can be made useful for code.

The intended downstream object is not a code generator. It is a scorer and
reranker. Given a before-state, an edit instruction or execution input, and a
set of candidate after-states, \modelname{} estimates which candidate best
matches the learned transition. The transition contract is:
\begin{equation}
  \hat z_{\mathrm{after}} =
  P(E(s_{\mathrm{before}}), A(a)).
\end{equation}

This paper reports the full evidence arc rather than only the strongest slice.
The early commit-message-conditioned edit substrate failed the intended
action-use gate. A later execution-substrate pivot produced non-collapsed
checkpoints and useful diagnostics. The final v0.9/v1.0 downstream replay then
found a narrow HumanEval WS-D win but did not clear the aggregate downstream
claim because MBPP-Plus was saturated by controls. This is an honest end state
for the artifact: reproducible systems evidence, negative model-quality
evidence on action use, one narrow positive reranking slice, and explicit
blocked claims.

\paragraph{Contributions.}
\begin{enumerate}[leftmargin=1.5em]
  \item A manifest-backed code world-model harness for dataset construction,
  training, retrieval, surprise, latent probes, downstream reranking, and
  paper-demo replay.
  \item A negative action-use result on commit-edit transitions, including
  v0.2 no-action dominance and unsupported latent semantic structure.
  \item A diagnostic execution-substrate result showing non-collapse,
  retrieval lift over controls, and mutation-surprise separation.
  \item A final downstream replay showing HumanEval WS-D lift over no-action on
  two seeds while closing the aggregate claim on MBPP-Plus saturation and
  broader representation/semantic gates.
\end{enumerate}

\section{Related Work}

\paragraph{Joint-embedding prediction.}
I-JEPA and V-JEPA learn feature predictions without pixel reconstruction
\citep{assran2023selfsupervisedlearningimagesjointembedding,bardes2024revisitingfeaturepredictionlearning}.
LeWorldModel studies stable end-to-end JEPA-style world modeling from pixels
\citep{maes2026leworldmodelstableendtoendjointembedding}. Anti-collapse
regularization is central to this family, with related ideas in VICReg, Barlow
Twins, and LeJEPA
\citep{bardes2022vicregvarianceinvariancecovarianceregularizationselfsupervised,zbontar2021barlowtwinsselfsupervisedlearning,balestriero2025lejepaprovablescalableselfsupervised}.

\paragraph{Code modeling and execution signals.}
Pretrained code models such as CodeBERT, GraphCodeBERT, UniXcoder, and CodeT5
learn program representations or generation models from code and natural
language
\citep{feng2020codebertpretrainedmodelprogramming,guo2021graphcodebertpretrainingcoderepresentations,guo2022unixcoderunifiedcrossmodalpretraining,wang2021codet5identifierawareunifiedpretrained}.
HumanEval, MBPP, APPS, CodeNet, and EvalPlus-style benchmark extensions provide
candidate tasks and labels for coding evaluation
\citep{chen2021evaluatinglargelanguagemodels,austin2021programsynthesislargelanguage,hendrycks2021measuringcodingchallengecompetence,puri2021codenetlargescaleaicode,liu2023codegeneratedchatgptreally}.
Generated tests have also been used for reranking code generations
\citep{chen2022codetcodegenerationgenerated}. Execution-aware pretraining and
execution-trace work provide a complementary signal for program behavior
\citep{zaremba2015learningexecute,ding2023tracedexecutionawarepretrainingsource,armengolestape2025iexecuteiunderstand}.
\modelname{} differs by treating code changes and execution records as
manifested latent transitions and by refusing to execute untrusted candidate
code during scoring.

\section{Model and Artifact Contract}

\modelname{} uses bounded code or execution states, an action encoder, a
transition predictor, and anti-collapse regularization. The exact module and
objective variants changed across research iterations, but the public contract
did not: every publishable run writes schema-versioned reports, artifact
manifests, checksums, and secret-scan evidence. Candidate code and generated
patches are untrusted. They may be parsed or scored as text, but they are not
imported, executed, or test-run by the model-scoring path.

The final paper-demo contract is deterministic. It replays checked-in v0.9 WS-D
completion labels and score rows for four slices: seed 42 and seed 1729 crossed
with HumanEval WS-D and MBPP-Plus WS-D. Replay mode records
\code{score\_source=replay\_existing\_scores}; it is not a fresh checkpoint
scoring run and not a live OpenRouter demo. The committed artifact id is
\code{demo\_report-e6fc06c328eed245}.

\section{Evidence Arc}

Table~\ref{tab:arc} summarizes the evidence by milestone. The final claim audit
is \artifact{docs/benchmark/V1_0_FINAL_CLAIM_AUDIT_2026-06-08.md}.

\begin{table}[t]
\centering
\caption{Evidence arc and public claim status.}
\label{tab:arc}
\small
\begin{tabular}{p{0.12\linewidth}p{0.25\linewidth}p{0.22\linewidth}p{0.30\linewidth}}
\toprule
Milestone & Main evidence & Open claim & Closed or blocked claim \\
\midrule
v0.2 & action-swap A10G run & systems pipeline & action-use, latent axes, downstream utility \\
v0.6 & execution substrate & diagnostics & downstream HumanEval/MBPP utility \\
v0.8 & execution repair & HumanEval WS-D slice & MBPP, magnitude probe, broad semantic gate \\
v0.9 & two-seed final gate suite & HumanEval WS-D slice & MBPP, representation, broad semantic gate \\
v1.0 & paper-demo replay & reproducible demo package & aggregate downstream usefulness \\
\bottomrule
\end{tabular}
\end{table}

\subsection{Negative action-use evidence}

The v0.2 action-swap/inverse-action run directly tested whether action
conditioning beat no-action on held-out retrieval. It did not. Text-action
Recall@1 was 0.263 and MRR was 0.370, while no-action reached Recall@1 0.441
and MRR 0.533. Exact-same-before, near-before, same-file, action-cluster,
edit-shape, and random action-contrast pools also failed to show the required
margin over no-action. Latent probes were explicitly unsupported for semantic
structure: lexical or metadata controls beat the latent views on available
targets.

\subsection{Execution-substrate diagnostics}

The execution substrate changed the state/action/next-state signal from noisy
commit messages to deterministic execution triples. This improved training
health and diagnostics. By v0.9, retrieval lift over controls passed on both
seeds: Recall@1 was 0.2538 and 0.2731, while all listed retrieval controls were
below 0.024 Recall@1. Mutation surprise AUC was 1.0 on both seeds. These are
training-health and diagnostic results; they do not by themselves prove
downstream coding usefulness.

Representation claims remain closed. The v0.9 \code{passed} probe reached
0.6552 and 0.6158 for \(\hat z_{\mathrm{after}}\), but no-action or before-state
controls were stronger. The \code{output\_magnitude\_bucket} target was
evaluable in v0.9 but tied a direct \code{z\_after} control. Broad semantic
surprise also remains closed because the semantic-decoy categories have zero
scorable aligned pairs in the final pack.

\section{Final Downstream Replay}

The final paper-demo replay is the paper's central downstream table. It is
derived from \artifact{docs/benchmark/v1_0/paper_demo/reports/paper_demo_report.json}
and summarized in Table~\ref{tab:demo}.

\begin{table}[t]
\centering
\caption{Final v1.0 paper-demo replay. The aggregate claim gate is closed even
though the HumanEval WS-D slices are open.}
\label{tab:demo}
\begin{tabular}{rlrrrrrl}
\toprule
Seed & Benchmark & CodeLeWM & No-action & LLM-order & Lexical & Lift & Gate \\
\midrule
42 & HumanEval WS-D & 0.9787 & 0.8723 & 0.1489 & 0.6596 & +10.64 & open \\
42 & MBPP-Plus WS-D & 1.0000 & 1.0000 & 0.1765 & 1.0000 & +0.00 & closed \\
1729 & HumanEval WS-D & 0.9787 & 0.8936 & 0.1489 & 0.6596 & +8.51 & open \\
1729 & MBPP-Plus WS-D & 1.0000 & 1.0000 & 0.1765 & 1.0000 & +0.00 & closed \\
\bottomrule
\end{tabular}
\end{table}

HumanEval WS-D is the narrow positive slice: CodeLeWM reaches pass@1 0.9787 on
both seeds and beats no-action by +10.64 and +8.51 points. The 95\% bootstrap
confidence intervals for lift over no-action are [2.13, 21.28] and
[2.13, 17.02], respectively. LLM-order is much weaker at pass@1 0.1489.

MBPP-Plus is the counterweight. CodeLeWM reaches pass@1 1.0000, but so do
no-action and lexical controls. The measured lift over no-action is exactly
+0.00 points on both seeds with confidence interval [0.00, 0.00]. This makes
the aggregate downstream claim closed. The evidence supports a narrow
HumanEval WS-D result, not a claim that \modelname{} generally improves patch
ranking.

\section{Why No-Action and Generalization Remain Hard}

The no-action baseline is strong because many code transitions are sparse:
large portions of the after-state are identical to the before-state, and many
candidate completions can be ranked using surface similarity. A world model
must learn when the action changes behavior, not merely when text resembles the
prompt. The v0.2 failure shows that a commit message is often too weak or too
underspecified to force this distinction.

The MBPP-Plus result exposes a different problem. The final WS-D MBPP-Plus pack
is saturated. CodeLeWM, no-action, and lexical all choose passing candidates at
rank 1 on both seeds. Saturation prevents a lift claim because the model did
not add measurable value beyond simple controls. This is not a bug in the paper
demo; it is the reason the aggregate claim gate must close. A future positive
claim needs harder, sufficiently large, public-safe downstream labels where
no-action and lexical controls do not already solve the slice.

The representation and semantic-surprise blockers reinforce the same point.
Healthy retrieval and mutation diagnostics show that the model is not collapsed
and can separate some generated decoys. They do not show stable semantic axes
or broad behavioral understanding. Those claims require probes and decoy packs
that beat controls across seeds and have enough scorable coverage.

\section{Claim Boundary}

Table~\ref{tab:claims} is the paper-level claim audit. It mirrors
\artifact{docs/papers/codelewm_final_claim_audit.md}.

\begin{table}[t]
\centering
\caption{Allowed and blocked public claims.}
\label{tab:claims}
\small
\begin{tabular}{p{0.28\linewidth}p{0.14\linewidth}p{0.47\linewidth}}
\toprule
Claim & Verdict & Evidence \\
\midrule
Reproducible artifact pipeline & allowed & manifests, scans, HF artifacts, paper demo \\
HumanEval WS-D reranking slice & allowed & 0.9787 pass@1 on both seeds \\
Broad coding improvement & blocked & aggregate downstream gate closed \\
MBPP-Plus/generalization & blocked & zero lift over no-action and lexical \\
Action-use model-quality & blocked & v0.2 no-action dominance \\
Semantic latent axes & blocked & probes do not beat controls \\
Standalone \code{p\_pass} downstream score & not recorded & final rows mark \code{p\_pass} missing \\
\bottomrule
\end{tabular}
\end{table}

\section{Reproducibility}

The source repository is \url{https://github.com/AbdelStark/CodeLeWM}. The
main artifact reports used in this paper are:

\begin{itemize}[leftmargin=1.5em]
  \item \artifact{docs/benchmark/V0_2_ACTION_SWAP_HF_RESULTS_2026-05-20.md}
  \item \artifact{docs/benchmark/EXECUTION_V0_8_RESULTS_2026-06-05.md}
  \item \artifact{docs/benchmark/EXECUTION_V0_9_RESULTS_2026-06-07.md}
  \item \artifact{docs/benchmark/PAPER_DEMO_V1_0_ARTIFACTS_2026-06-08.md}
  \item \artifact{docs/benchmark/V1_0_FINAL_CLAIM_AUDIT_2026-06-08.md}
  \item \artifact{docs/benchmark/v1_0/paper_demo/manifest.json}
\end{itemize}

The deterministic replay command for the final demo is:
\begin{verbatim}
uv run scripts/paper-demo --out docs/benchmark/v1_0/paper_demo --overwrite --json
\end{verbatim}
The paper build command is:
\begin{verbatim}
scripts/build-codelewm-final-paper
\end{verbatim}

\section{Limitations}

\begin{itemize}[leftmargin=1.5em]
  \item The final positive result is a slice result, not a broad benchmark
  result.
  \item The final paper demo replays checked-in v0.9 score rows. It is
  reproducible and manifest-backed, but it is not a fresh checkpoint scoring
  run.
  \item MBPP-Plus saturation prevents measuring lift on that slice.
  \item The representation probes do not support naming semantic dimensions.
  \item Broad semantic surprise remains blocked by scorable-pair coverage.
  \item Candidate code remains untrusted; the model-scoring path does not
  execute candidate patches.
  \item The results do not compare against large pretrained code models as
  generators. \modelname{} is a scorer/reranker artifact.
\end{itemize}

\section{Conclusion}

\modelname{} is a completed reproducible research artifact, not a broad coding
assistant. The artifact discipline worked: data, checkpoints, reports, claim
gates, manifests, and scans can be regenerated and inspected. The original
action-use hypothesis did not work: no-action stayed stronger on commit-edit
retrieval. The execution substrate gave healthier model diagnostics and a
narrow HumanEval WS-D downstream win, but MBPP-Plus saturation and representation
blockers keep the public model-quality claim closed. The most useful conclusion
is methodological: code world-model claims should be published with their
baselines and failure gates intact, especially when no-action and lexical
controls are strong.

\bibliographystyle{plainnat}
\bibliography{two_substrate_references}

\end{document}
